% !TeX TXS-program:compile = txs:///pdflatex/[-shell-escape]
\documentclass[11pt]{beamer}
\usepackage[utf8]{inputenc}
\usepackage[T1,T2A]{fontenc}
\usepackage[russian]{babel}
\usepackage{color}
\usepackage{listings}
\usepackage{calc}
\usepackage{graphicx}
\usepackage{epstopdf}
\usepackage{hyperref}
\hypersetup{unicode,colorlinks}
\usepackage{csquotes}
\usepackage{upquote}
\usepackage{cprotect}
\usetheme[progressbar=head,numbering=fraction,block=fill]{metropolis}
\usepackage{dejavu}
\usepackage{etoolbox}
\usepackage{bm}
\usepackage{minted}
%\usepackage[defaultsans]{droidsans}
%\usepackage[defaultmono]{droidmono}
%\makeatletter
%\input droid/t1fdm.fd
%\makeatother

\makeatletter
\DeclareRobustCommand*{\bfseries}{%
    \not@math@alphabet\bfseries\mathbf
    \fontseries\bfdefault\selectfont
    \boldmath
}
\makeatother

\makeatletter
\chardef\straightquote@code=\catcode`'
\chardef\backquote@code=\catcode``
\catcode`'=\active \catcode``=\active
\patchcmd{\@noligs}
{\textasciigrave}
{\fixedtextasciigrave}
{}{}
\newcommand{\fixedtextasciigrave}{%
    \makebox[.5em]{\fontencoding{TS1}\fontfamily{fvs}\selectfont\textasciigrave}% Vera Sans
}
\catcode\lq\'=\straightquote@code
\catcode\lq\`=\backquote@code
\makeatletter

\lstset{
    language=Haskell,
    inputencoding=utf8,
    extendedchars=true,
    breaklines=true,
    escapeinside=!!,
    tabsize=4,
    breakatwhitespace=true,
    keepspaces=true,
    basicstyle=\ttfamily,
    commentstyle=\color{gray}\itshape,
    keywordstyle=
    % keywordstyle=\color{blue}, % TODO https://tex.stackexchange.com/questions/415777/avoid-highlighting-keywords-following-certain-words-in-listings
    % addkeywords={Type}
    % deletekeywords={Monad, String, return, array, show},
    % otherkeywords={::, ->}
}
\lstset{
    literate={а}{{\selectfont\char224}}1
    {б}{{\selectfont\char225}}1
    {в}{{\selectfont\char226}}1
    {г}{{\selectfont\char227}}1
    {д}{{\selectfont\char228}}1
    {е}{{\selectfont\char229}}1
    {ё}{{\"e}}1
    {ж}{{\selectfont\char230}}1
    {з}{{\selectfont\char231}}1
    {и}{{\selectfont\char232}}1
    {й}{{\selectfont\char233}}1
    {к}{{\selectfont\char234}}1
    {л}{{\selectfont\char235}}1
    {м}{{\selectfont\char236}}1
    {н}{{\selectfont\char237}}1
    {о}{{\selectfont\char238}}1
    {п}{{\selectfont\char239}}1
    {р}{{\selectfont\char240}}1
    {с}{{\selectfont\char241}}1
    {т}{{\selectfont\char242}}1
    {у}{{\selectfont\char243}}1
    {ф}{{\selectfont\char244}}1
    {х}{{\selectfont\char245}}1
    {ц}{{\selectfont\char246}}1
    {ч}{{\selectfont\char247}}1
    {ш}{{\selectfont\char248}}1
    {щ}{{\selectfont\char249}}1
    {ъ}{{\selectfont\char250}}1
    {ы}{{\selectfont\char251}}1
    {ь}{{\selectfont\char252}}1
    {э}{{\selectfont\char253}}1
    {ю}{{\selectfont\char254}}1
    {я}{{\selectfont\char255}}1
    {А}{{\selectfont\char192}}1
    {Б}{{\selectfont\char193}}1
    {В}{{\selectfont\char194}}1
    {Г}{{\selectfont\char195}}1
    {Д}{{\selectfont\char196}}1
    {Е}{{\selectfont\char197}}1
    {Ё}{{\"E}}1
    {Ж}{{\selectfont\char198}}1
    {З}{{\selectfont\char199}}1
    {И}{{\selectfont\char200}}1
    {Й}{{\selectfont\char201}}1
    {К}{{\selectfont\char202}}1
    {Л}{{\selectfont\char203}}1
    {М}{{\selectfont\char204}}1
    {Н}{{\selectfont\char205}}1
    {О}{{\selectfont\char206}}1
    {П}{{\selectfont\char207}}1
    {Р}{{\selectfont\char208}}1
    {С}{{\selectfont\char209}}1
    {Т}{{\selectfont\char210}}1
    {У}{{\selectfont\char211}}1
    {Ф}{{\selectfont\char212}}1
    {Х}{{\selectfont\char213}}1
    {Ц}{{\selectfont\char214}}1
    {Ч}{{\selectfont\char215}}1
    {Ш}{{\selectfont\char216}}1
    {Щ}{{\selectfont\char217}}1
    {Ъ}{{\selectfont\char218}}1
    {Ы}{{\selectfont\char219}}1
    {Ь}{{\selectfont\char220}}1
    {Э}{{\selectfont\char221}}1
    {Ю}{{\selectfont\char222}}1
    {Я}{{\selectfont\char223}}1
    {`}{{{\fixedtextasciigrave}}}1
}

% \usemintedstyle{solarized-light}
\newminted[haskell]{haskell}{
    escapeinside=``,
    beameroverlays=true,
    texcomments=true,
    autogobble=true,
    fontsize=\small
}
\newmintinline[haskinline]{haskell}{
    escapeinside=``,
    beameroverlays=true,
    texcomments=true
}

\author{Алексей Романов}
\subtitle{Функциональное программирование на Haskell}
%\logo{}
\institute{МИЭТ}
\subject{Функциональное программирование на Haskell}
%\setbeamercovered{transparent}
%\setbeamertemplate{navigation symbols}{}

\usepackage{tikz}
\title{Лекция 15: комбинаторы парсеров}

\begin{document}
\begin{frame}[plain]
  \maketitle
\end{frame}

\begin{frame}[fragile]
  \frametitle{От текста к программе}
  \begin{itemize}
    \item В лабораторной 2 выражения задавались значениями типа \haskinline|Expression|. Как получить их из текста?
      \begin{columns}[c,onlytextwidth]
        \begin{column}{0.52\linewidth}
          \begin{haskellsmall}
            "2 + 3 * 4"
            Op (Val 2) Plus
               (Op (Val 3) Times (Val 4))
          \end{haskellsmall}
        \end{column}
        \begin{column}{0.46\linewidth}
          \centering
          \begin{tikzpicture}[
            every node/.style={font=\footnotesize\ttfamily},
            edge from parent/.style={draw,semithick},
            level distance=0.75cm,
            level 1/.style={sibling distance=1.7cm},
            level 2/.style={sibling distance=1.3cm}]
            \node {Op: Plus}
              child {node {Val: 2}}
              child {node {Op: Times}
                child {node {Val: 3}}
                child {node {Val: 4}}};
          \end{tikzpicture}
        \end{column}
      \end{columns}
      \pause
    \item \emph{Синтаксический анализ} (parsing) превращает текст в структуру, например дерево выражения.
    \item Вычисление результата по этому дереву "--- отдельная задача: для SimpleLang её решает \haskinline|eval|.
  \end{itemize}
\end{frame}

\begin{frame}[fragile]
  \frametitle{Тип парсера}
  \begin{haskell}
    newtype Parser a = Parser
      { runParser :: String -> Maybe (a, String) }
  \end{haskell}
  \begin{itemize}
    \item \haskinline|a| "--- тип результата. Успех возвращает результат и \emph{остаток входа}, неудача "--- \haskinline|Nothing|.
      \pause
    \item Можно использовать также \hackage{text} (\haskinline|Text|) и \hackage{bytestring} (\haskinline|ByteString|). Для простоты берём \haskinline|String|.
    \item В Megaparsec эти варианты объединяет класс \href{https://hackage.haskell.org/package/megaparsec/docs/Text-Megaparsec-Stream.html}{\haskinline{Stream}}. \haskinline|Text| содержит символы, \haskinline|ByteString| "--- байты.
  \end{itemize}
\end{frame}

\begin{frame}[fragile]
  \frametitle{Один символ}
  Парсер \haskinline|satisfy| читает и возвращает первый символ, если он удовлетворяет предикату. Если символ не подходит или вход пуст, разбор завершается неудачей.
  \pause
  \begin{haskellsmall}
    satisfy :: (Char -> Bool) -> Parser Char
    satisfy predicate =!\pause! Parser $ \input ->
      case input of
        c : rest | predicate c -> Just (c, rest)
        _                      -> Nothing
  \end{haskellsmall}
\end{frame}

\begin{frame}[fragile]
  \frametitle{Заданный символ}
  Парсер \haskinline|char| читает и возвращает заданный символ. Если первый символ другой или вход пуст, разбор завершается неудачей.
  \pause
  \begin{haskellsmall}
    char :: Char -> Parser Char
    char c =!\pause! satisfy (== c)
  \end{haskellsmall}
  \begin{haskellsmall}
    ghci> runParser (char 'a') "abc" !\pause!
    Just ('a',"bc")
    ghci> runParser (char 'a') "xyz" !\pause!
    Nothing
  \end{haskellsmall}
\end{frame}

\begin{frame}[fragile]
  \frametitle{Парсер как функтор}
  \begin{itemize}
    \item Изменяем результат, сохраняя остаток входа и неудачу:
      \begin{haskellsmall}
        instance Functor Parser where
          fmap g (Parser parse) =!\pause! Parser $ \input ->
            fmap (\(value, rest) -> (g value, rest))
                 (parse input)
      \end{haskellsmall}
    \item Здесь внутренний \haskinline|fmap| работает с \haskinline|Maybe|.
  \end{itemize}
\end{frame}

\begin{frame}[fragile]
  \frametitle{Цифра и её значение}
  \begin{itemize}
    \item \haskinline|digit| читает одну десятичную цифру и возвращает её как символ.
      \pause
      \begin{haskellsmall}
        digit :: Parser Char
        digit =!\pause! satisfy (\c -> '0' <= c && c <= '9')
      \end{haskellsmall}
      \pause
    \item \haskinline|digitValue| читает цифру и возвращает её числовое значение.
      \pause
      \begin{haskellsmall}
        digitValue :: Parser Int
        digitValue =!\pause! (\c -> fromEnum c - fromEnum '0')
                     <$> digit
      \end{haskellsmall}
  \end{itemize}
\end{frame}

\begin{frame}[fragile]
  \frametitle{Последовательность парсеров: Applicative}
  \begin{itemize}
    \item Второй парсер получает остаток входа после первого.
      \begin{haskellsmall}
        instance Applicative Parser where
          pure value = Parser $ \input ->
            Just (value, input)
          Parser parseG <*> Parser parseX =
            Parser $ \input -> do
              (g, rest1) <- parseG input
              (value, rest2) <- parseX rest1
              pure (g value, rest2)
      \end{haskellsmall}
      \pause
    \item Здесь \haskinline|do| работает в \haskinline|Maybe|. При неудаче любой части получаем \haskinline|Nothing|.
    \item Упражнение: определите \haskinline|liftA2| напрямую через \haskinline|runParser|.
  \end{itemize}
\end{frame}

\begin{frame}[fragile]
  \frametitle{Собираем результат}
  \haskinline|twoDigits| читает две цифры подряд и возвращает пару их числовых значений.
  \pause
  \begin{haskell}
    twoDigits :: Parser (Int, Int)
    twoDigits =!\pause! liftA2 (,) digitValue digitValue
  \end{haskell}
  \begin{haskellsmall}
    ghci> runParser twoDigits "42!" !\pause!
    Just ((4,2),"!")
  \end{haskellsmall}
\end{frame}

\begin{frame}[fragile]
  \frametitle{Парсер в скобках}
  \haskinline|parens| запускает переданный парсер между круглыми скобками. Скобки читаются, но в результат входит только значение внутреннего парсера.
  \pause
  \begin{haskell}
    parens :: Parser a -> Parser a
    parens p =!\pause! char '(' *> p <* char ')'
  \end{haskell}
  \begin{haskellsmall}
    ghci> runParser (parens twoDigits) "(42)!" !\pause!
    Just ((4,2),"!")
    ghci> runParser (parens twoDigits) "(4x)!" !\pause!
    Nothing
  \end{haskellsmall}
\end{frame}

\begin{frame}[fragile]
  \frametitle{Заданная строка}
  Парсер \haskinline|string| читает и возвращает заданную строку. Если она не совпадает с началом входа, разбор завершается неудачей.
  \pause
  \begin{haskellsmall}
    string :: String -> Parser String
    string =!\pause! traverse char
  \end{haskellsmall}
  \pause
  Для двух символов определение \haskinline|traverse| раскрывается так:
  \begin{haskellsmall}
    traverse char "ab" !\pause!
      !$=$! (:) <$> char 'a'
            <*> ((:) <$> char 'b' <*> pure [])
  \end{haskellsmall}
  \begin{haskellsmall}
    ghci> runParser (string "ab") "abc" !\pause!
    Just ("ab","c")
    ghci> runParser (string "ab") "ac!" !\pause!
    Nothing
  \end{haskellsmall}
\end{frame}

\begin{frame}[fragile]
  \frametitle{Выбор: Alternative}
  \begin{itemize}
    \item Класс из \href{https://hackage.haskell.org/package/base/docs/Control-Applicative.html}{\haskinline{Control.Applicative}} добавляет неудачу и выбор:
      \begin{haskellsmall}
        class Applicative f => Alternative f where
          empty :: f a
          (<|>) :: f a -> f a -> f a
      \end{haskellsmall}
    \item Это основные методы; повторение обсудим отдельно.
      \pause
    \item Для нашего парсера сначала пробуем левый вариант. При неудаче пробуем правый \emph{на исходном входе}:
      \begin{haskellsmall}
        instance Alternative Parser where
          empty = Parser $ const Nothing
          Parser p <|> Parser q = Parser $ \input ->
            p input <|> q input
      \end{haskellsmall}
    \item Справа \haskinline{(<|>)} работает с \haskinline|Maybe|.
  \end{itemize}
\end{frame}

\begin{frame}[fragile]
  \frametitle{Возврат к началу альтернативы}
  \haskinline|abOrAc| разбирает и возвращает строку \haskinline|"ab"| или \haskinline|"ac"|.
  \pause
  \begin{haskellsmall}
    abOrAc :: Parser String
    abOrAc =!\pause! string "ab" <|> string "ac"
  \end{haskellsmall}
  \begin{haskellsmall}
    ghci> runParser abOrAc "ac!" !\pause!
    Just ("ac","!")
  \end{haskellsmall}
  \begin{itemize}
    \item После несовпадения с \haskinline|"ab"| правая ветка пробует \haskinline|"ac"| на исходном входе.
      \pause
    \item Успех слева завершает выбор. Позднейшая неудача другого парсера не заставляет пересмотреть этот выбор.
  \end{itemize}
\end{frame}

\begin{frame}[fragile]
  \frametitle{Успешный выбор не пересматривается}
  Парсер \haskinline|p| выбирает между \haskinline|"a"| и \haskinline|"ab"|, а затем читает \haskinline|"c"|.
  \pause
  \begin{haskellsmall}
    p :: Parser String
    p =!\pause! (string "a" <|> string "ab") <* string "c"
  \end{haskellsmall}
  \begin{haskellsmall}
    ghci> runParser p "abc" !\pause!
    Nothing
  \end{haskellsmall}
  \begin{itemize}
    \item Левая ветка успешно читает \haskinline|"a"|. Остаток "--- \haskinline|"bc"|.
    \item Следующий \haskinline|string "c"| не срабатывает. К варианту \haskinline|"ab"| внутри завершившегося выбора уже не возвращаемся.
  \end{itemize}
\end{frame}

\begin{frame}[fragile]
  \frametitle{Границы альтернативы}
  Включим чтение конечного \haskinline|"c"| в каждую ветку. Парсер \haskinline|q| выбирает между \haskinline|"ac"| и \haskinline|"abc"|.
  \pause
  \begin{haskellsmall}
    q :: Parser String
    q =!\pause! (string "a" <* string "c")
        <|> (string "ab" <* string "c")
  \end{haskellsmall}
  \begin{haskellsmall}
    ghci> runParser q "abc" !\pause!
    Just ("ab","")
  \end{haskellsmall}
  \begin{itemize}
    \item Неудача на \haskinline|"c"| теперь происходит внутри левой ветки. Правая ветка запускается на исходном входе.
  \end{itemize}
\end{frame}

\begin{frame}[fragile]
  \frametitle{Повторение: many и some}
  \haskinline|some| повторяет парсер хотя бы один раз, \haskinline|many| допускает ноль повторений. Результаты собираются в список.
  \pause
  \begin{haskellsmall}
    many :: Alternative f => f a -> f [a]
    some :: Alternative f => f a -> f [a]

    some p =!\pause! (:) <$> p <*> many p
    many p =!\pause! some p <|> pure []
  \end{haskellsmall}
  Это взаимные определения по умолчанию в \haskinline|Alternative|.
  \begin{haskellsmall}
    ghci> runParser (some digit) "123x" !\pause!
    Just ("123","x")
    ghci> runParser (many digit) "xyz" !\pause!
    Just ("","xyz")
  \end{haskellsmall}
\end{frame}

\begin{frame}[fragile]
  \frametitle{Повторяемый парсер должен продвигаться}
  \begin{itemize}
    \item Что произойдёт при разборе?
      \begin{haskell}
        runParser (many (pure 'x')) "abc"
      \end{haskell}
      \pause
    \item \haskinline|pure 'x'| всегда успешно завершается, но не читает вход. Рекурсия в \haskinline|many| не заканчивается.
    \item В нашей реализации это зацикливание. Некоторые библиотеки диагностируют такую ошибку.
    \item Не передавайте в \haskinline|many| или \haskinline|some| парсер, который может успешно завершиться без чтения входа.
  \end{itemize}
\end{frame}

\begin{frame}[fragile]
  \frametitle{Лексический и синтаксический анализ}
  \begin{itemize}
    \item \emph{Лексический анализ} (lexing) выделяет токены: числа, имена, знаки операций и т.д. Пробелы и комментарии обычно пропускаются.
    \item Например, \haskinline|"12 + total"| содержит три токена: \haskinline|12|, \haskinline|+| и \haskinline|total|.
      \pause
    \item \emph{Синтаксический анализ} (parsing) объединяет токены по правилам грамматики, строя дерево выражения или программы.
      \pause
    \item У нас оба этапа совмещены: парсеры читают символы, а комбинаторы собирают дерево.
    \item Можно сначала получить список токенов, а затем разобрать этот список по грамматике.
  \end{itemize}
\end{frame}

\begin{frame}[fragile]
  \frametitle{Числа и пробелы}
  \begin{itemize}
    \item \haskinline|natural| читает непустую последовательность цифр и возвращает неотрицательное целое число.
      \pause
      \begin{haskellsmall}
        natural :: Parser Integer
        natural =!\pause! read <$> some digit
      \end{haskellsmall}
      \pause
    \item \haskinline|spaces| читает ноль или больше пробельных символов.
      \pause
      \begin{haskellsmall}
        spaces :: Parser String
        spaces =!\pause! many (satisfy isSpace)
      \end{haskellsmall}
    \item \haskinline|isSpace| импортируем из \href{https://hackage.haskell.org/package/base/docs/Data-Char.html}{\haskinline{Data.Char}}.
  \end{itemize}
\end{frame}

\begin{frame}[fragile]
  \frametitle{Пробелы после токена}
  \begin{itemize}
    \item \haskinline|lexeme| запускает парсер, затем читает пробелы. Возвращает результат переданного парсера.
      \pause
      \begin{haskellsmall}
        lexeme :: Parser a -> Parser a
        lexeme p =!\pause! p <* spaces
      \end{haskellsmall}
      \pause
    \item \haskinline|symbol| читает заданную строку и следующие пробелы.
      \pause
      \begin{haskellsmall}
        symbol :: String -> Parser String
        symbol =!\pause! lexeme . string
      \end{haskellsmall}
    \item Чтобы разрешить начальные пробелы, начните итоговый парсер с \haskinline|spaces *>|.
  \end{itemize}
\end{frame}

\begin{frame}[fragile]
  \frametitle{Разделители и скобки}
  \haskinline|between| запускает парсер между открывающим и закрывающим разделителями, возвращая только внутренний результат.
  \pause
  \begin{haskellsmall}
    between :: Parser open -> Parser close ->
               Parser a -> Parser a
    between open close p =!\pause! open *> p <* close
  \end{haskellsmall}
  \pause
  \haskinline|sepBy| разбирает список элементов с разделителями между ними. Пустой список допустим. Разделитель после последнего элемента остаётся во входе.
  \pause
  \begin{haskellsmall}
    sepBy :: Parser a -> Parser sep -> Parser [a]
    sepBy p sep = !\pause!
      ((:) <$> p <*> many (sep *> p)) <|> pure []
  \end{haskellsmall}
\end{frame}

\begin{frame}[fragile]
  \frametitle{Список чисел}
  \haskinline|numbers| разбирает список неотрицательных целых чисел в квадратных скобках, разделённых запятыми. Между токенами допускаются пробелы.
  \pause
  \begin{haskellsmall}
    numbers :: Parser [Integer]
    numbers =!\pause! between (symbol "[") (symbol "]")
      (sepBy (lexeme natural) (symbol ","))
  \end{haskellsmall}
  \begin{itemize}
    \item Упражнение: что останется от \haskinline|"[1, 2, 3]!"|? А что будет с \haskinline|"[1,]"|?
  \end{itemize}
\end{frame}

\begin{frame}[fragile]
  \frametitle{Нужно ли разобрать весь вход?}
  \haskinline|eof| успешно завершается только на пустом остатке входа и возвращает \haskinline|()|.
  \pause
  \begin{haskellsmall}
    eof :: Parser ()
    eof =!\pause! Parser $ \input -> case input of
      [] -> Just ((), [])
      _  -> Nothing
  \end{haskellsmall}
  \pause
  \haskinline|parseAll| запускает парсер, требуя полного разбора входа. Возвращает только результат.
  \pause
  \begin{haskellsmall}
    parseAll :: Parser a -> String -> Maybe a
    parseAll p input = !\pause!
      fst <$> runParser (p <* eof) input
  \end{haskellsmall}
  \pause
  \begin{haskellsmall}
    ghci> parseAll natural "12xyz" !\pause!
    Nothing
  \end{haskellsmall}
\end{frame}

\begin{frame}[fragile]
  \frametitle{Парсер как монада}
  \begin{haskellsmall}
    instance Monad Parser where
      Parser parse >>= next = Parser $ \input -> do
        (value, rest) <- parse input
        runParser (next value) rest
  \end{haskellsmall}
  \begin{itemize}
    \item \haskinline|next| по результату первого парсера выбирает следующий.
    \item Остаток входа передаётся неявно, как состояние в \haskinline|State|. Неудача передаётся через \haskinline|Maybe|.
      \pause
    \item Большую часть обычной грамматики удобно выразить через \haskinline|Applicative| и \haskinline|Alternative|. \haskinline|do| также можно использовать для читаемости.
  \end{itemize}
\end{frame}

\begin{frame}[fragile]
  \frametitle{Когда следующая часть зависит от результата}
  \begin{itemize}
    \item \haskinline|counted| читает неотрицательное целое число, двоеточие и указанное число символов. Возвращает эти символы, например \haskinline|"3:abc"| даёт \haskinline|"abc"|.
      \pause
      \begin{haskell}
        counted :: Parser String
        counted =!\pause! do
          (n :: Int) <- read <$> some digit
          char ':'
          replicateM n (satisfy (const True))
      \end{haskell}
      \pause
    \item Число запусков парсера зависит от прочитанного \haskinline|n|. Здесь используем \haskinline|Monad|.
    \item \haskinline|replicateM| "--- из \href{https://hackage.haskell.org/package/base/docs/Control-Monad.html}{\haskinline{Control.Monad}}.
    \item Упражнение: предскажите остаток для \haskinline|"2:abc"| и результат для \haskinline|"3:ab"|.
  \end{itemize}
\end{frame}

\begin{frame}[fragile]
  \frametitle{Дерево выражения}
  \begin{itemize}
    \item Часть типов из SimpleLang (лабораторная 2):
      \begin{haskell}
        data Expression
          = Var String
          | Val Integer
          | Op Expression Bop Expression
          deriving (Show, Eq)

        data Bop = Plus | Minus | Times
          deriving (Show, Eq)
      \end{haskell}
    \item Пока разбираем числа, переменные, скобки и \haskinline|+|. Унарного минуса нет.
    \item Имена состоят из латинских букв.
  \end{itemize}
\end{frame}

\begin{frame}[fragile]
  \frametitle{Первая грамматика: только сложение}
  \begin{haskell}
    expr0 = atom { '+' atom }
    atom  = natural | name | '(' expr0 ')'
  \end{haskell}
  \begin{itemize}
    \item Это \emph{грамматика} в нотации \href{https://en.wikipedia.org/wiki/Extended_Backus\%E2\%80\%93Naur_form}{EBNF}, расширении \href{https://en.wikipedia.org/wiki/Backus\%E2\%80\%93Naur_form}{формы Бэкуса--Наура (BNF)}, а не Haskell. Фигурные скобки обозначают ноль или больше повторений, вертикальная черта "--- выбор.
      \pause
    \item \haskinline|atom| "--- неделимое выражение на этом уровне: число, переменная или выражение в скобках.
    \item Выражение в скобках снова содержит \haskinline|expr0|: грамматика рекурсивна.
  \end{itemize}
\end{frame}

\begin{frame}[fragile]
  \frametitle{Имя переменной}
  \haskinline|name| читает непустую последовательность латинских букв и следующие пробелы, возвращая имя.
  \pause
  \begin{haskellsmall}
    name :: Parser String
    name =!\pause! lexeme $ some $ satisfy $ \c ->
      ('a' <= c && c <= 'z') ||
      ('A' <= c && c <= 'Z')
  \end{haskellsmall}
\end{frame}

\begin{frame}[fragile]
  \frametitle{Парсер атома}
  \haskinline|atom| разбирает число, переменную или выражение в круглых скобках и возвращает соответствующее дерево.
  \pause
  \begin{haskellsmall}
    atom :: Parser Expression
    atom =!\pause! Val <$> lexeme natural
       <|> Var <$> name
       <|> between (symbol "(") (symbol ")") expr0
  \end{haskellsmall}
  \begin{itemize}
    \item Все три ветки возвращают \haskinline|Expression|.
    \item Перед рекурсивным вызовом читается открывающая скобка.
  \end{itemize}
\end{frame}

\begin{frame}[fragile]
  \frametitle{Парсер сложения}
  \haskinline|expr0| разбирает цепочку слагаемых и собирает дерево сложений слева направо.
  \pause
  \begin{haskellsmall}
    expr0 :: Parser Expression
    expr0 =!\pause! assemble <$> atom
                     <*> many (symbol "+" *> atom)
      where
        assemble first rest =
          foldl (\left right -> Op left Plus right)
                first rest
  \end{haskellsmall}
  \begin{haskellsmall}
    ghci> parseAll expr0 "a + 2 + 3" !\pause!
    Just (Op (Op (Var "a") Plus (Val 2))
             Plus (Val 3))
  \end{haskellsmall}
  \begin{itemize}
    \item Список хвостовых слагаемых превращаем в дерево слева направо.
  \end{itemize}
\end{frame}

\begin{frame}[fragile]
  \frametitle{Добавим умножение: что получится?}
  \haskinline|operator| читает знак сложения или умножения и возвращает соответствующий конструктор операции.
  \pause
  \begin{haskellsmall}
    operator :: Parser Bop
    operator =!\pause! Plus  <$ symbol "+"
           <|> Times <$ symbol "*"
  \end{haskellsmall}
  \begin{itemize}
    \item Попробуем разбирать любую цепочку \haskinline|atom| и \haskinline|operator| слева направо, как сложение.
    \item Какое дерево получится для \haskinline|"2 + 3 * 4"|?
      \pause
      \begin{haskellsmall}
        Op (Op (Val 2) Plus (Val 3)) Times (Val 4)
      \end{haskellsmall}
    \item Это \haskinline|(2 + 3) * 4|. Чтобы получить привычный результат, нужно задать \emph{приоритеты} операторов.
  \end{itemize}
\end{frame}

\begin{frame}[fragile]
  \frametitle{Приоритеты операторов}
  \begin{haskell}
    expr0 = expr1 { ('+' | '-') expr1 }
    expr1 = atom  { '*' atom }
    atom  = natural | name | '(' expr0 ')'
  \end{haskell}
  \begin{itemize}
    \item \haskinline|expr0| разбирает сложение и вычитание, \haskinline|expr1| "--- умножение.
    \item Умножение собирается в поддерево раньше, чем это поддерево становится аргументом сложения.
      \pause
    \item \haskinline|expr0|, \haskinline|expr1| "--- наши имена уровней. При добавлении операторов с другими приоритетами можно добавить новые уровни.
    \item Скобки позволяют снова использовать любой уровень внутри \haskinline|atom|.
  \end{itemize}
\end{frame}

\begin{frame}[fragile]
  \frametitle{Парсер бинарного оператора}
  \haskinline|binary| читает заданный знак и возвращает функцию построения узла дерева с выбранной операцией.
  \pause
  \begin{haskellsmall}
    binary :: Bop -> String ->
              Parser (Expression -> Expression -> Expression)
    binary op text =!\pause! (\left right -> Op left op right)
                  <$ symbol text
  \end{haskellsmall}
  \pause
  \haskinline|addOp| читает плюс или минус и возвращает соответствующую функцию построения узла.
  \pause
  \begin{haskellsmall}
    addOp :: Parser (Expression -> Expression -> Expression)
    addOp =!\pause! binary Plus "+" <|> binary Minus "-"
  \end{haskellsmall}
\end{frame}

\begin{frame}[fragile]
  \frametitle{Ассоциативность: как собирать цепочку?}
  \begin{itemize}
    \item При одинаковом приоритете: как разобрать \haskinline|8 - 3 + 1|?
      \pause
    \item \emph{Левая ассоциативность}: \haskinline|(8 - 3) + 1|. Собираем дерево через \haskinline|foldl|.
    \item \haskinline|chainLeft| разбирает цепочку операндов с операторами и собирает её слева направо.
      \pause
      \begin{haskellsmall}
        chainLeft :: Parser a ->
                     Parser (a -> a -> a) -> Parser a
        chainLeft p op =!\pause! assemble <$> p
                                 <*> many ((,) <$> op <*> p)
          where
            assemble = foldl (\left (g, right) ->
                               g left right)
      \end{haskellsmall}
  \end{itemize}
\end{frame}

\begin{frame}[fragile]
  \frametitle{Парсер с приоритетами}
  \begin{itemize}
    \item \haskinline|expr1| разбирает цепочку умножений атомов.
      \pause
      \begin{haskellsmall}
        expr1 :: Parser Expression
        expr1 =!\pause! chainLeft atom (binary Times "*")
      \end{haskellsmall}
      \pause
    \item \haskinline|expr0| разбирает сложения и вычитания выражений уровня \haskinline|expr1|.
      \pause
      \begin{haskellsmall}
        expr0 :: Parser Expression
        expr0 =!\pause! chainLeft expr1 addOp
      \end{haskellsmall}
    \item Используем прежний \haskinline|atom|. Его скобки вызывают новый \haskinline|expr0|.
    \item Для большего числа приоритетов добавляем уровни \haskinline|exprN|: операнды каждого уровня разбирает следующий, более высокий.
  \end{itemize}
\end{frame}

\begin{frame}[fragile]
  \frametitle{Приоритеты и скобки: примеры}
  \begin{haskellsmall}
    ghci> parseAll expr0 "2 + 3 * 4" !\pause!
    Just (Op (Val 2) Plus
             (Op (Val 3) Times (Val 4)))
    ghci> parseAll expr0 "(2 + 3) * 4" !\pause!
    Just (Op (Op (Val 2) Plus (Val 3))
             Times (Val 4))
    ghci> parseAll expr0 "8 - 3 + 1" !\pause!
    Just (Op (Op (Val 8) Minus (Val 3))
             Plus (Val 1))
  \end{haskellsmall}
\end{frame}

\begin{frame}[fragile]
  \frametitle{Почему не записать правило напрямую?}
  \begin{haskell}
    expr0 = expr0 '+' atom | atom
  \end{haskell}
  \begin{itemize}
    \item Это \emph{левая рекурсия}: правило вызывает себя до чтения входа.
      \pause
    \item Прямой перевод в наш парсер снова вызывает \haskinline|expr0| с тем же входом и не заканчивается.
    \item Переписанное правило читает первый атом, а затем повторяет хвост:
      \begin{haskell}
        expr0 = atom { '+' atom }
      \end{haskell}
    \item Грамматики описывают один язык, но второй вариант подходит для нашего способа разбора.
  \end{itemize}
\end{frame}

\begin{frame}[fragile]
  \frametitle{От учебного парсера к Megaparsec}
  \begin{itemize}
    \item \hackage{megaparsec}: тот же стиль композиции, но с позициями во входе и сообщениями об ошибках.
      \begin{haskellsmall}
        import Data.Void (Void)
        import Text.Megaparsec
        import Text.Megaparsec.Char
        import qualified Text.Megaparsec.Char.Lexer as L

        type Parser = Parsec Void String
      \end{haskellsmall}
    \item \haskinline|Void| из \href{https://hackage.haskell.org/package/base/docs/Data-Void.html}{\haskinline{Data.Void}} означает отсутствие наших собственных типов ошибок.
      \pause
    \item Это \emph{отдельный пример}: наш тип \haskinline|Parser| и одноимённые комбинаторы заменяем библиотечными.
  \end{itemize}
\end{frame}

\begin{frame}[fragile]
  \frametitle{Megaparsec: лексические комбинаторы}
  \begin{itemize}
    \item У нас: \haskinline|lexeme p = p <* spaces|.
    \item \href{https://hackage.haskell.org/package/megaparsec/docs/Text-Megaparsec-Char-Lexer.html#v:lexeme}{\haskinline{L.lexeme}} принимает парсер пропусков явно:
      \begin{haskellsmall}
        lexeme :: Parser a -> Parser a
        lexeme = L.lexeme space
      \end{haskellsmall}
      Аналогично: \href{https://hackage.haskell.org/package/megaparsec/docs/Text-Megaparsec-Char-Lexer.html#v:symbol}{\haskinline{symbol = L.symbol space}}. Здесь \href{https://hackage.haskell.org/package/megaparsec/docs/Text-Megaparsec-Char.html#v:space}{\haskinline{space}} пропускает пробельные символы.
    \item \href{https://hackage.haskell.org/package/megaparsec/docs/Text-Megaparsec-Char-Lexer.html#v:space}{\haskinline{L.space}} поддерживает комментарии: \href{https://hackage.haskell.org/package/megaparsec/docs/Text-Megaparsec-Char-Lexer.html#v:skipLineComment}{\haskinline{L.skipLineComment}} и \href{https://hackage.haskell.org/package/megaparsec/docs/Text-Megaparsec-Char-Lexer.html#v:skipBlockComment}{\haskinline{L.skipBlockComment}}.
    \item \href{https://hackage.haskell.org/package/megaparsec/docs/Text-Megaparsec-Char-Lexer.html#v:signed}{\haskinline{L.signed}} добавляет к числовому парсеру необязательный знак.
    \item Есть \href{https://hackage.haskell.org/package/megaparsec/docs/Text-Megaparsec-Char-Lexer.html#v:indentBlock}{комбинаторы для учёта отступов}.
  \end{itemize}
\end{frame}

\begin{frame}[fragile]
  \frametitle{Megaparsec: запуск и ошибка}
  \href{https://hackage.haskell.org/package/megaparsec/docs/Text-Megaparsec-Char-Lexer.html#v:decimal}{\haskinline{L.decimal}} читает неотрицательное целое число.
  \begin{haskellsmall}
    ghci> parseTest (L.decimal <* eof :: Parser Integer) "12x" !\pause!
    1:3:
      |
    1 | 12x
      |   ^
    unexpected 'x'
    expecting digit or end of input
  \end{haskellsmall}
  \begin{itemize}
    \item \href{https://hackage.haskell.org/package/megaparsec/docs/Text-Megaparsec.html}{\haskinline{parseTest}} печатает результат или ошибку. \haskinline|parse| возвращает \haskinline|Either| с ошибкой или результатом.
  \end{itemize}
\end{frame}

\begin{frame}[fragile]
  \frametitle{Megaparsec: выбор после чтения входа}
  \begin{itemize}
    \item В ветках из отдельных \haskinline|char| без \haskinline|try| на \haskinline|"ac"| первая ветка читает \haskinline|'a'| и завершается ошибкой. Правая ветка \emph{не запускается}.
      \pause
    \item \haskinline|try| позволяет откатить потребление при неудаче:
      \begin{haskellsmall}
        abOrAc :: Parser Char
        abOrAc = try (char 'a' *> char 'b')
             <|> (char 'a' *> char 'c')
      \end{haskellsmall}
    \item Здесь лучше вынести общий префикс: \haskinline{char 'a' *> (char 'b' <|> char 'c')}.
    \item Библиотечный \haskinline|string| уже откатывает несовпадение целой строки. Это не относится к произвольной последовательности парсеров.
  \end{itemize}
\end{frame}

\begin{frame}[fragile]
  \frametitle{Как тестировать парсер}
  \begin{itemize}
    \item Примеры: пустой вход, вложенные скобки, пробелы, неверный символ, незакрытая скобка, лишний текст после выражения.
    \item Приоритеты и ассоциативность проверяем по \emph{дереву}, а не только по результату вычисления.
  \end{itemize}
\end{frame}

\begin{frame}[fragile]
  \frametitle{Свойство обратного разбора}
  \begin{itemize}
    \item Печать выражения и обратный разбор возвращают исходное дерево:
      \begin{haskellsmall}
        parseAll (spaces *> expr0)
          (renderExpression expression)
          === Just expression
      \end{haskellsmall}
    \item \haskinline|renderExpression| из лабораторной 2 пишет выражение в нашем синтаксисе, например с полными скобками. Производный \haskinline|Show| печатает конструкторы и здесь не подходит.
    \item Генератор должен создавать только выражения, поддерживаемые парсером: например, без отрицательных констант, с допустимыми именами.
  \end{itemize}
\end{frame}

\begin{frame}[fragile]
  \frametitle{Упражнения}
  \begin{itemize}
    \item Добавьте унарный минус. Решите, какой у него приоритет и допускается ли несколько минусов подряд.
    \item Добавьте правоассоциативное возведение в степень: \haskinline|2 ^ 3 ^ 2| должно означать \haskinline|2 ^ (3 ^ 2)|. Расширьте тип дерева.
    \item Проверьте свойство обратного разбора для \haskinline|renderExpression| из лабораторной 2 с Hedgehog.
    \item Необязательно: реализуйте парсер для \haskinline|DietStatement| или \haskinline|Statement|. Запустите результат через \haskinline|runSimpler| или \haskinline|run| соответственно.
  \end{itemize}
\end{frame}

\begin{frame}[fragile]
  \frametitle{Комбинаторы парсеров в других языках}
  \begin{itemize}
    \item Java: \href{https://github.com/jparsec/jparsec}{jparsec} и \href{https://github.com/google/mug/blob/master/dot-parse/README.md}{DotParse} "--- композиция парсеров средствами Java.
    \item Scala: \href{https://com-lihaoyi.github.io/fastparse/}{FastParse} "--- последовательность, выбор и повторение в виде операторов и методов.
    \item Rust: \href{https://github.com/rust-bakery/nom}{nom} "--- текстовые и бинарные форматы, в том числе чтение по частям.
    \item C++: \href{https://www.boost.org/doc/libs/latest/libs/spirit/doc/x3/html/index.html}{Boost.Spirit X3} и \href{https://github.com/boostorg/spirit_x4}{Spirit X4} "--- грамматика как выражение C++, собирающее парсеры.
    \item Python: \href{https://parsy.readthedocs.io/en/latest/overview.html}{parsy} "--- комбинаторы в духе Parsec.
  \end{itemize}
\end{frame}

\begin{frame}[fragile]
  \frametitle{Другие способы разбора}
  \begin{itemize}
    \item \emph{Регулярные выражения}: часто используются для распознавания отдельных токенов, например чисел и имён.
    \item \emph{Рекурсивный спуск вручную}: функции для правил грамматики явно управляют позицией во входе и выбором веток.
    \item \emph{Генераторы парсеров}: по грамматике в отдельном файле создают код парсера. Примеры: \href{https://haskell-happy.readthedocs.io/en/latest/}{Happy} для Haskell и \href{https://www.antlr.org/}{ANTLR} для нескольких языков.
    \item \emph{Инкрементальный разбор}: обновляет дерево после изменения текста. \href{https://tree-sitter.github.io/tree-sitter/}{Tree-sitter} генерирует такие парсеры для использования в редакторах.
  \end{itemize}
\end{frame}

\begin{frame}[fragile]
  \frametitle{Дополнительное чтение: комбинаторы парсеров}
  \begin{itemize}
    \item \href{https://www.cs.nott.ac.uk/~pszgmh/monparsing.pdf}{Hutton, Meijer: Monadic Parser Combinators} "--- учебная реализация.
    \item \href{https://markkarpov.com/tutorial/megaparsec.html}{Mark Karpov: Megaparsec tutorial} "--- практическое руководство.
    \item \href{https://dl.acm.org/doi/10.1145/3471874.3472984}{Willis, Wu: Design Patterns for Parser Combinators (Functional Pearl)} "--- организация парсеров.
    \item \href{https://hackage.haskell.org/package/megaparsec/docs/Text-Megaparsec.html}{Документация Megaparsec}.
    \item \hackage{parser-combinators} "--- готовые комбинаторы, в том числе \href{https://hackage.haskell.org/package/parser-combinators/docs/Control-Monad-Combinators-Expr.html}{\haskinline{makeExprParser}} для разбора выражений с заданными операторами и уровнями приоритета.
    \item \href{https://hackage.haskell.org/package/base/docs/Control-Applicative.html}{\haskinline{Control.Applicative}} "--- общие комбинаторы.
  \end{itemize}
\end{frame}
\begin{frame}[fragile]
  \frametitle{Дополнительное чтение: разбор в целом}
  \begin{itemize}
    \item \href{https://craftinginterpreters.com/contents.html}{Robert Nystrom: Crafting Interpreters} "--- бесплатная онлайн-книга. Главы 4--6: лексический анализ, деревья и рекурсивный спуск на Java; глава 17: разбор по Пратту.
    \item \href{https://dthain.github.io/books/compiler/}{Douglas Thain: Introduction to Compilers and Language Design} "--- бесплатный учебник. Главы 3--5: лексический анализ, синтаксический анализ и его реализация.
    \item Обе книги посвящены реализации языков программирования в целом, а не только синтаксическому анализу.
  \end{itemize}
\end{frame}
\end{document}
